% OCaml Workshop 2026 — extended abstract (~2 pages).
% Fill the bracketed prompts IN YOUR OWN WORDS. The CFP desk-rejects proposals
% "largely written by LLMs"; the comments are notes on WHAT to say, not text to submit.
% Deadline July 1 AoE; submit at https://types-hotcrp.paris.inria.fr/ocaml26/
% Section headings are optional (the accepted one-pagers often omit them) — delete the
% \section*{} lines for a flowing WasML-style version if you prefer.

\documentclass[11pt,a4paper]{article}
\usepackage[margin=2.4cm]{geometry}
\usepackage[T1]{fontenc}
\usepackage{microtype}
\usepackage[hidelinks]{hyperref}

\title{\vspace{-2.5em}\texttt{ocaml-mmtk}: OCaml on the Memory Management Toolkit
       as a testbed for functional-language GC}
\author{[Presenting author] \and [co-authors] \\[2pt]
        \small [affiliation(s)]}
\date{}

\begin{document}
\maketitle
\thispagestyle{empty}

% ===== 1. Intro =====================================================
% Your validated ~120-word abstract, lightly expanded if you like:
% MMTk -> we present ocaml-mmtk (MMTk = OCaml 5.5's only GC, plans swap at startup;
% self-hosts) -> first functional, statically-typed language as an MMTk substrate
% (vs OpenJDK/Julia/Ruby) -> the two questions (immutability -> read-barrier-free? ;
% mutation-rate-dominated cost -> generalises to Haskell/Erlang?).
\textit{[Intro paragraph — your prose.]}

% ===== 2. Background: OCaml's GC and the read-barrier question =======
\section*{Background}
% Why OCaml is the natural test case. OCaml 5's own collector (ICFP'20\footnote{...}):
%  - per-domain copying MINOR heap -> promotes into a SHARED, non-moving,
%    mostly-concurrent MARK-AND-SWEEP major;
%  - an SATB (snapshot-at-the-beginning) write barrier, and -- deliberately -- NO READ
%    BARRIER: they chose the STW ParMinor over a concurrent minor specifically to keep
%    the C API's raw-read assumption (C stubs read OCaml values directly);
%  - the authors lean on immutability explicitly: the initialising writes (the only
%    ones for immutable data) need no barrier, and reads need no barrier.
% So the read-barrier-free question is not arbitrary -- it is the design axis OCaml's
% own authors reasoned about. Read: RESEARCH_QUESTIONS.md baseline section.
\textit{[Background paragraph — your prose.]}

% ===== 3. ocaml-mmtk (the artifact) ================================
\section*{\texttt{ocaml-mmtk}}
% High-level, NO build minutiae:
%  - we implemented MMTk's VMBinding for OCaml: object model, PRECISE root scanning
%    (reusing caml_do_roots / caml_scan_stack), stop-the-world, write barriers;
%  - MMTk is the ONLY collector, for bytecode AND native code; native works via TLAB
%    nursery-aliasing (OCaml's inlined bump allocator points at an MMTk block) with NO
%    code-generator change; plans swap at startup; the compiler self-hosts.
% RQ4 contrast (one or two sentences -- the "GC-friendly language" story):
%  - OCaml's precise rooting + no naked pointers meant NONE of the conservative-pinning
%    machinery the Julia and CRuby MMTk reports\footnote{...}\footnote{...} needed;
%  - but the COORDINATION impedance persisted -- e.g. a stop-the-world handshake counter
%    underflow once let the GC scan domains that had not actually stopped (root-caused
%    under rr) -- a generalisable GC-runtime-interface finding.
\textit{[Artifact paragraph — your prose.]}

% ===== 4. Early results ============================================
\section*{Early results}
% Preliminary, SEQUENTIAL (state the caveat -- do NOT claim parallel wins):
%  - throughput vs stock OCaml 5.5.0 (native, memory parity): the default generational
%    plan is at parity-or-better on most benchmarks; a read-barrier-free concurrent-
%    marking plan runs allocation-heavy code up to ~1.9x faster;
%  - that plan's SATB write barrier is essentially free (<0.1% on the most write-heavy
%    bench), and moving marking off the pause cuts the max stop-the-world pause ~3-4x;
%  - removing MMTk's redundant allocation-time zeroing (safe by OCaml's own init-before-
%    GC discipline) recovered ~15-22% on allocation-heavy code, no pause cost.
% CAVEAT sentence: single-heap-size numbers are misleading; the rigorous heap-size-
% normalised comparison vs stock is in progress. Read: README quick panel (reproducible),
% RESEARCH_QUESTIONS.md RQ1 native result + RQ8.
\textit{[Early-results paragraph / short list — your prose.]}

% ===== 5. + 6. Research questions and outlook ======================
\section*{Research questions and outlook}
% RQ1 (flagship): does the RESIDUAL throughput cost of read-barrier-free collection scale
% with a program's MUTATION RATE / lifetime dispersion (Dolan'25\footnote{...}) rather
% than its allocation rate? -- regress the gap across the mutability spectrum
% (pure-functional <-> ref/Bytes-heavy). If so, it generalises to the functional family
% (Haskell, Erlang). Next: an LXR-style reference-counting plan.
\textit{[RQ1 paragraph — your prose.]}
%
% RQ7 / the synthesis (the testbed payoff, and the close): the data already splits the
% win -- the generational plan wins throughput (cheap nursery), the concurrent-marking
% plan wins latency (off-pause marking). OCaml's OWN collector is generational +
% concurrent-marking + non-moving-major + SATB + read-barrier-free -- so the natural next
% plan, "Bactrian", re-expresses OCaml's own GC inside MMTk: simultaneously the cleanest
% MMTk-vs-vanilla comparison (isolating framework overhead from collector design) and the
% low-latency vehicle RQ1 needs. The talk presents the platform, this early evidence, and
% this direction.
\textit{[RQ7 + closing paragraph — your prose.]}

% ===== Footnotes / references ======================================
% Fill the \footnote{...} markers above, or use a short bibliography:
%  ICFP'20  Sivaramakrishnan et al., Retrofitting Parallelism onto OCaml, ICFP 2020.
%  LXR      Zhao, Blackburn & McKinley, Low-Latency High-Throughput GC, PLDI 2022.
%  Julia    de Souza Amorim et al., Reconsidering GC in Julia, ISMM 2025.
%  CRuby    Wang et al., Reworking Memory Management in CRuby, ISMM 2025.
%  Dolan25  Dolan, Lifetime Dispersion and Generational GC, ISMM 2025.
%  MMTk     Lin, Blackburn, Hosking & Norrish, Rust as a Language for ... GC, ISMM 2016.

\end{document}
